\paragraph{Direct fine-tuning results.}
We train three students for 1,000 updates in each run. Across 512 generated
anchors, 351 pass the graph checks and 349 have at least one allowed force-shifted
candidate. Composition-level mean ESS is about 1.27--1.87 out of 8. The weights
are therefore concentrated on few candidates, limiting resolution of the local
target with this pool size.

\begin{table}[t]
\centering
\caption{Graph-valid outputs after direct physical fine-tuning, out of 640
attempts per run. The frozen model receives no updates; all other variants have
the same additional training budget.}
\begin{tabular}{lrrrr}
\toprule
Run & Frozen & Replay & Unweighted escort & Complete work\\
\midrule
1 &418&340&330&344\\
2 &377&345&371&362\\
\bottomrule
\end{tabular}\label{tab:local-work-direct}
\end{table}

Among paired outputs that are both graph-valid, complete-work training gives
lower energy than unweighted-escort training by 0.00758 eV/atom (paired 95\%
interval $[-0.01424,-0.00113]$ for the signed difference). Force RMS is lower by
0.16174 eV/\AA\ ($[-0.27660,-0.05590]$). Both runs favor complete work, and
the all-output comparison has the same direction.

However, all three directly fine-tuned students perform worse than the frozen
model. Complete work has 55.16\% graph validity versus 62.11\% for the frozen
model, and energy on paired, jointly valid outputs is \emph{higher} by
0.10587 eV/atom ($[0.09009,0.12168]$). Replay also worsens performance after
training on the mixed targets. These results motivate subtracting the matched
replay update rather than using the physical student directly.

Teacher preparation uses 6,232 raw energy/force queries: 702 at anchors and
5,530 at allowed candidates. The unweighted escort could omit candidate-energy
queries, and replay could omit all physical queries. Scoring the outputs of the
four models uses a further 10,240 queries, separate from generation cost.
